\section{RDFS and OWL}

\subsection{From facts to knowledge}

RDF triples express \emph{facts}, but facts alone are not always \emph{knowledge}.
The difference is \textbf{inference}: deriving statements that were never explicitly
written---for example, knowing Alice is an \texttt{Employee} because she
\texttt{worksFor} TechCorp, even if nobody typed \texttt{a :Employee}.

RDF Schema (RDFS) and the Web Ontology Language (OWL) are the semantic layers that
turn raw RDF into structured knowledge that both machines and LLMs can reason over
(\href{https://bryon.io/from-facts-to-knowledge-the-layer-cake-of-rdfs-and-owl-e84819d8075d}{Bryon Jacob, 2025}).

\subsection{The Semantic Web layer cake}

Tim Berners-Lee described the Semantic Web as a stack of layers, each building on
the one below
(\href{https://bryon.io/from-facts-to-knowledge-the-layer-cake-of-rdfs-and-owl-e84819d8075d}{Bryon Jacob, 2025}):

\begin{table}[htbp]
  \centering
  \small
  \caption{Semantic Web layer cake (bottom to top)}
  \label{tab:layer-cake}
  \renewcommand{\arraystretch}{1.2}
  \begin{tabularx}{\textwidth}{@{}l X@{}}
    \toprule
    \textbf{Layer} & \textbf{What it adds} \\
    \midrule
    Raw data & SQL, CSV, JSON, spreadsheets---messy enterprise silos \\
    RDF & A common graph of triples: ``Alice works at TechCorp'' \\
    RDFS & Types and relationships: ``If X worksFor Y, X is a Person and Y is an Organization'' \\
    OWL & Logic: transitivity, equivalence, disjointness, cardinality constraints \\
    Reasoners & Deterministic inference and consistency checking \\
    LLM interface & Natural language in; explanations out \\
    \bottomrule
  \end{tabularx}
\end{table}

Each step adds power. Raw data says \texttt{SELECT * FROM emp\_tbl WHERE dept\_id = 42};
RDF says ``Alice works at TechCorp''; RDFS adds typing rules; OWL concludes
``If you work at a company in a city, you work in that city.''

\subsection{RDF Schema (RDFS)}

RDFS provides a basic vocabulary for describing \textbf{classes} (types of things)
and \textbf{properties} (how things relate).

\paragraph{Classes and properties}

\begin{verbatim}
# Define what things ARE
:Person a rdfs:Class ;
    rdfs:label "Person" ;
    rdfs:comment "A human being" .

:Organization a rdfs:Class ;
    rdfs:label "Organization" .

# Define how things RELATE
:worksFor a rdf:Property ;
    rdfs:domain :Person ;
    rdfs:range :Organization ;
    rdfs:label "works for" .
\end{verbatim}

\paragraph{Domain and range (type inference)}

Given a single fact:

\begin{verbatim}
tc:employee-alice-johnson :worksFor tc:org-techcorp .
\end{verbatim}

RDFS inference yields:

\begin{verbatim}
tc:employee-alice-johnson a :Person .
tc:org-techcorp a :Organization .
\end{verbatim}

The LLM does not need to guess Alice's type from naming conventions---the schema
states it.

\paragraph{Class hierarchies}

\begin{verbatim}
:Employee rdfs:subClassOf :Person .
:Manager rdfs:subClassOf :Employee .
:SeniorEngineer rdfs:subClassOf :Employee .

tc:employee-alice-johnson a :SeniorEngineer .
# Inferred: also :Employee and :Person
\end{verbatim}

Hierarchies let systems answer ``What kind of thing is Alice?'' at the right level
of specificity.

\subsection{OWL: logic on top of data}

If RDFS is arithmetic, OWL is calculus
(\href{https://bryon.io/from-facts-to-knowledge-the-layer-cake-of-rdfs-and-owl-e84819d8075d}{Bryon Jacob, 2025}).
OWL adds logical constructs for sophisticated, decidable reasoning (based on Description Logic).

\paragraph{Equivalence and disjointness}

\begin{verbatim}
:Company owl:equivalentClass :Corporation .

:Person owl:disjointWith :Organization .
# Something cannot be both a Person and an Organization
\end{verbatim}

\paragraph{Transitive properties}

\begin{verbatim}
:locatedIn a owl:TransitiveProperty .

techcorp:hq :locatedIn city:seattle .
city:seattle :locatedIn state:washington .
state:washington :locatedIn country:usa .

# Inferred:
# techcorp:hq :locatedIn state:washington .
# techcorp:hq :locatedIn country:usa .
\end{verbatim}

Without OWL, traversing location hierarchies requires custom query logic. With OWL,
a reasoner computes the closure deterministically.

\paragraph{Inverse properties}

\begin{verbatim}
:employs owl:inverseOf :worksFor .

:alice :worksFor :techcorp .
# Inferred: :techcorp :employs :alice .
\end{verbatim}

\paragraph{Property chains}

\begin{verbatim}
:workedOnBy owl:propertyChainAxiom ( :managedBy :worksFor ) .

tc:project-kg :managedBy tc:employee-alice-johnson .
tc:employee-alice-johnson :worksFor tc:org-techcorp .
# Inferred: tc:project-kg :workedOnBy tc:org-techcorp .
\end{verbatim}

\paragraph{Consistency and cardinality}

OWL reasoners catch logical errors before they reach downstream systems:

\begin{verbatim}
:something a :Person, :Organization .   # inconsistent (disjoint)

:Project rdfs:subClassOf [
    a owl:Restriction ;
    owl:onProperty :hasManager ;
    owl:cardinality 1
] .
:project123 :hasManager :alice, :bob .  # too many managers
\end{verbatim}

\subsection{Why LLMs benefit from RDFS and OWL}

LLMs read OWL axioms almost as easily as English---\texttt{:locatedIn a
owl:TransitiveProperty} is unambiguous. The breakthrough is pairing LLMs with
\textbf{deterministic reasoners} that either prove a conclusion or do not
(\href{https://bryon.io/from-facts-to-knowledge-the-layer-cake-of-rdfs-and-owl-e84819d8075d}{Bryon Jacob, 2025}):

\begin{itemize}
  \item \textbf{Automatic type understanding:} reasoners compute that Alice is a
        \texttt{Person}, \texttt{Employee}, and \texttt{SeniorEngineer}; the LLM
        queries proven facts, not guesses.
  \item \textbf{Consistency guarantees:} contradictions are flagged before answers
        are generated.
  \item \textbf{Transitive reasoning:} ``Does Alice work in the USA?'' becomes a
        graph traversal computed by a reasoner, not a multi-hop hallucination risk.
  \item \textbf{Query expansion:} subclass relationships are materialized
        automatically.
\end{itemize}

The LLM orchestrates: understand intent, translate to formal queries, call
reasoning tools, explain results in natural language.

\subsection{The 80/20 rule of semantic modeling}

You do not need to model everything. A small RDFS/OWL core often delivers most of
the value (\href{https://bryon.io/from-facts-to-knowledge-the-layer-cake-of-rdfs-and-owl-e84819d8075d}{Bryon Jacob, 2025}):

\begin{table}[htbp]
  \centering
  \small
  \caption{High-value RDFS and OWL constructs}
  \label{tab:rdfs-owl-constructs}
  \renewcommand{\arraystretch}{1.15}
  \begin{tabularx}{\textwidth}{@{}l X@{}}
    \toprule
    \textbf{RDFS (start here)} & \textbf{OWL (add selectively)} \\
    \midrule
    \texttt{rdfs:subClassOf} & \texttt{owl:TransitiveProperty} \\
    \texttt{rdfs:domain} & \texttt{owl:inverseOf} \\
    \texttt{rdfs:range} & \texttt{owl:disjointWith} \\
    \texttt{rdfs:subPropertyOf} & \texttt{owl:equivalentClass} \\
    \bottomrule
  \end{tabularx}
\end{table}

As James Hendler noted: ``A little semantics goes a long way''
(\href{https://bryon.io/from-facts-to-knowledge-the-layer-cake-of-rdfs-and-owl-e84819d8075d}{Bryon Jacob, 2025}).

\subsection{Practical patterns for LLM integration}

\paragraph{Semantic validation}

Define valid business objects with OWL, then let a reasoner validate incoming
data. The LLM translates constraint violations into plain language (e.g.\ ``This
order is invalid because it has no items'').

\paragraph{Inference-powered question answering}

Combine natural-language understanding with deterministic inference---for
example, defining \texttt{hasAncestor} as transitive over \texttt{hasParent}, so
``Who are Alice's ancestors?'' is answered by a reasoner, not by the LLM inventing
family trees.
